\documentclass[12pt,a4paper,oneside]{book}


\usepackage[utf8]{inputenc}
\usepackage[a4paper,inner=3.5cm,outer=2.5cm]{geometry}

\usepackage[titletoc,title,toc,page]{appendix}
\usepackage{verbatim}
\usepackage{placeins}
\usepackage{listings}
\usepackage{hyperref}
\usepackage[italian]{babel}
\usepackage{tikz}
\usepackage{parskip}
\usepackage{url}

\usepackage{graphicx}
\usepackage{blindtext}
\usepackage{chngcntr}
\counterwithin{table}{chapter}

\usepackage{newlfont}
\usepackage{fancyhdr}
\usepackage{indentfirst}
\usepackage[utf8]{inputenc}
\usepackage{float}
\usepackage{hyperref}
\usepackage[capitalize,noabbrev]{cleveref}
\usepackage{soul}
\usepackage[font=footnotesize,labelfont=bf]{caption}

\usepackage{multirow}
\usepackage{hyphenat}
\hyphenation{mate-mati-ca recu-perare}

\usepackage{lscape} 

\usepackage{natbib}
\bibliographystyle{alpha}
\setcitestyle{super,open={[},close={]}}
\usepackage{hyperref}
\newcommand{\rom}[1]{\uppercase\expandafter{\romannumeral #1\relax}}

\usepackage{pdfpages}

\begin{document}
% Per spostare i vari elementi più su o più giù gioca con i valori di vspace che ci sono tra uno e l'altro
\pagestyle{empty}
\newgeometry{
    left=20mm,
    right=20mm,
    top=20mm,
    bottom=20mm
}

\begin{titlepage}

\begin{center}

% marchio di ateneo
\includegraphics[width=6.5cm,height=4.7cm]{img/marchio-di-ateneo.png}

\vspace{15mm}


% \Large is 14.4pt
{\Large{\bf{Tesina Storia e Politiche del digitale}}}

\vspace{15mm}

{\Huge{\bf I costi materiali delle }}\\
\vspace{3mm}
{\Huge{\bf tecnologie del XXIesimo secolo}}\\
\vspace{3mm}
\end{center}

\vspace{20mm}

\hfill
\begin{minipage}[t]{0.40\textwidth}\raggedleft
{\Large{\bf Presentata da: \\ Alessandro Nanni \\ mat. 
 0001027757}}
\end{minipage}

\vspace{90mm}

\rule[0.5cm]{15.8cm}{0.6mm}

\begin{center}
{\large{\bf Anno Accademico 2024/2025\\}}
\end{center}

\end{titlepage}

\restoregeometry

\pagenumbering{gobble}


\topmargin=-1cm
\tableofcontents

\pagenumbering{arabic}
\raggedbottom


\chapter{Le terre rare}
\pagestyle{plain}
\setcounter{page}{1}

\section{Scoperta e utilizzo}
La scoperta delle terre rare risale al 1789, quando una «dura pietra nera» fu trovata nella miniera della città svedese di Ytterby. Il chimico Johan Gadolin non fu in grado di identificare la pietra, la quale chiamò «terra rara». La ricerca e isolamento di nuovi materiali portò alla scoperta di 17 elementi metallici tra il Diciottesimo e Ventesimo secolo: le terre rare, o in inglese, \textit{rare earth elements} (REE).
Nonostante questi metalli siano stati scoperti abbastanza presto, la loro utilità divenne apparente solamente dopo la Seconda Guerra Mondiale, in seguito alla diffusione delle televisioni. Negli anni 70 si iniziò a sfruttare le capacità magnetiche delle terre rare per creare supermagneti, fondamentali per creare motori elettrici.

Al giorno d'oggi, le REE sono indispensabili per tecnologie di comunicazione e informazione, grazie alle loro proprietà di semiconduttori. Ci sono due tipi di terre rare, quelle leggere e più comuni (LREE), e quelle più pesanti, rare e difficili da estrarre (HREE). Le terre rare leggere sono presenti nella crosta terrestre in quantità simili a quelle del rame, fondamentali per le \textit{green technologies} che stanno venendo lentamente adottare per ridurre l'impatto umano sull'ambiente. Anche quelle pesanti, pur essendo meno comuni, sono comunque molto più diffuse rispetto a metalli preziosi come l'oro.
Data l'abbondanza e importanza di questi metalli, un nome più appropriato potrebbe essere Terre Vitali?

\section{Miniere e geopolitica}
L'amplia presenza di terre rare nel suolo non implica la possibilità di realizzare un supermagnete o una batteria con facilità. Anzi, l'intero processo che va dall'estrazione alla produzione finale è complesso, costoso e richiede tecnologie avanzate, oltre a porre importanti sfide ambientali e geopolitiche.

La Cina e gli Stati Uniti sono stati i primi a gettarsi nel mercato delle terre rare, attorno al 1950. L'infrastruttura per l'estrazione cinese si è particolarmente sviluppata tra il 1980 e 1990, grazie a riforme che hanno promosso lo sviluppo tecnologico e industriale, avendo già notato l'utilità nel breve e nel lungo termine degli REE. In quegli anni gli Stati Uniti invece hanno assunto un ruolo marginale nell'estrazione ed esportazione di terre rare, probabilmente a causa di limitazioni legali nella gestione di materiale radioattivo.

Al giorno d'oggi la Cina estrae, raffina e vende l'80\% delle terre rare. Questa posizione di dominio non è dovuta solamente alla lungimiranza del governo centrale e alla disposizione degli strumenti e ricchezze necessarie: lavoro in nero a prezzo basso, standard ambientali bassi, miniere illegali e mercati neri hanno portato a prezzi con cui l'occidente non è in grado di competere. 

Il 70-80\% delle terre rare è estratta in miniere illegali a Bayan-Obo, o nelle montagne di Jiangxi, nella Cina meridionale. Qua uomini e donne vengono pagati poche centinaia di euro al mese per scavare e schiacciare rocce rosse in quello che sembra un gigantesco formicaio, sempre attivo.
«Like a human anthill, the mountain was mined twenty-four hours a day, seven days a week» \citep[p. 24]{pitron2020rare}.
L'estrazione di queste rocce è stata vietata dalle autorità cinesi nel 2016, e i minatori illegali hanno dovuto pagare pesanti multe. Inoltre, enormi scorte di metalli promesse a mercati esteri sono state sequestrate al porto di Canton, e dozzine di trafficanti furono incarcerati.
Nonostante ciò alcuni minatori si sono chiusi nei punti più inospitali delle montagne, dove si arricchiscono di nascosto. I materiali che scavano poi vengono messi in vendita in un mercato nero che li raffina e poi esporta in tutto il mondo.

Vengono impiegati principalmente due metodi per estrarre le terre rare, entrambi presentano un forte impatto ambientale con un conseguente effetto negativo sulla salute dei minatori dovuto all'impiego di sostanze chimiche tossiche. Il primo consiste nello scavare lo strato superficiale del terreno e creare vasche di estrazione, dove si versano reagenti chimici per separare i metalli dalla terra. Queste vasche, se non propriamente isolate, possono contaminare le falde acquifere e compromettere sistemi idrici.
Il secondo metodo consiste nel trivellare il terreno e inserire tubi in PVC, attraverso i quali vengono inserite sostanze chimiche direttamente nel sottosuolo. Anche in questo caso si crea una vasca di estrazione, con i medesimi effetti negativi del primo metodo. Inoltre, spesso questi tubi in PVC vengono abbandonati nel sito dello scavo.

Gli effetti negativi dello scavo di REE non sono limitati al sottosuolo terrestre; coinvolgono anche

\begin{itemize}
\item fiumi
\begin{quote}
\small
In 2006, some sixty companies producing indium — a rare metal used in the manufacture of certain solar panel technologies — released tonnes of chemicals into the Xiang River in Hunan, jeopardising the meridional province’s drinking water and the health of its residents. \citep[p. 25]{pitron2020rare}
\end{quote}

\item coltivazioni ed ecosistemi
\begin{quote}
\small
'Men and women, wearing no more than basic face masks, work in areas thick with black particles and acid fumes. It’s hell.' To complete the picture are toxic pits of chemical discharges from the plants, fields of poisoned corn, acid rain, and more. \citep[p. 26]{pitron2020rare}
\end{quote}

\item cittadini
\begin{quote}
\small
Rare earths have cost the community dearly. The hair of young men barely thirty years of age has suddenly turned white. Children grow up without developing any teeth. In 2010, the Chinese press reported that sixty-six Dalahai residents had died of cancer. \citep[p. 29]{pitron2020rare}
\end{quote}
    
\end{itemize}

\section{Interesse di Trump}
Il monopolio cinese sulle terre rare potrebbe essere uno dei motivi che ha recentemente spinto il presidente degli Stati Uniti, Donald Trump, a rivolgere l'attenzione verso paesi ricchi di questi materiali strategici, in particolare la Groenlandia e l'Ucraina.

Trump sostiene che il controllo sulla Groenlandia fornirebbe non solo maggiore ricchezza e indipendenza per gli Stati Uniti, ma anche un notevole vantaggio militare, installando missili balistici e radar sul fondo oceanico per rilevare navi russe e cinesi.
Gli Stati Uniti hanno colto l'importanza strategica della Groenlandia fin dal secondo dopoguerra, quando nel 1946 il presidente Truman provò a comprare l'isola dai danesi. In data 28 marzo 2025, il vicepresidente JD Vance ha visitato la Groenlandia, affermando che il governo locale ha «investito troppo poco nelle persone della Groenlandia e trascurato anche l'architettura della sicurezza di questa terra straordinaria, meravigliosa e popolata da persone eccezionali» \cite{Phi}. 
Le risorse minerarie dell'isola risultano ancora in gran parte inesplorate, dato che molte di esse sono sepolte sotto uno spesso strato di ghiaccio, e si stima che siano di più di quelle presenti in Cina. Costruire miniere per estrarre terre rare potrebbe essere la chiave degli Stati Uniti per non dover più importare REE dalla Cina.
Tuttavia in Groenlandia sono presenti numerose leggi ambientali e sociali che renderebbero la costruzione di nuove miniere un processo lungo e costoso. Inoltre sono già state emanate leggi che vietano l'estrazione di gas e petrolio presente nel fondo oceanico attorno all'isola.

Un mese prima della visita del vicepresidente Vance in Groenlandia, Trump ha annunciato, nel corso di una conferenza stampa, la conclusione di un accordo con il presidente ucraino Volodymyr Zelenskyy riguardante un investimento congiunto nelle risorse minerarie presenti sul territorio ucraino. Questa era la seconda iterazione di un accordo dove inizialmente erano proposti all'Ucraina aiuti militari equivalenti a \$500mld, che Zelenskyy aveva rapidamente rifiutato sostenendo «Non firmerò qualcosa che dieci generazioni di ucraini dovranno ripagare» \cite{guardian2025trumpdeal}.
Il nuovo accordo prevede un fondo di investimento per la ricostruzione congiunto tra Stati Uniti e Ucraina. L'Ucraina contribuirà al fondo con il 50\% di tutti i ricavi ottenuti dalla futura monetizzazione delle risorse naturali di proprietà del governo ucraino.
Questo accordo è poi fallito, dato che il suolo ucraino non era ricco di REE tanto quanto l'amministazione americana credeva. Anche se ci fossero terre rare da estrarre, e ignorando le difficoltà legate al conflitto con la Russia, ci vorrebbero comunque almeno 15 anni prima di avere miniere attive in grado di produrre ricchezze.\\
«Se cercate minerali critici, l'Ucraina non è certo il posto giusto dove trovarli» – dichiara Jack Lifton, presidente esecutivo del \textit{Critical Minerals Institute}. «È una fantasia. Tutto questo non ha senso. C'è evidentemente un altro tipo di agenda dietro. Faccio fatica a credere che qualcuno a Washington possa davvero pensare che abbia senso procurarsi terre rare in Ucraina» \cite{ieee2023ukraine}.

\section{Riciclo e recupero}
Attualmente il riciclo delle terre rare è limitato da una raccolta inefficace, alti costi di smontaggio e dall’assenza di metodi economicamente sostenibili per il recupero.
La piccola quantità che viene effettivamente riciclata proviene da:
\begin{itemize}
    \item magneti permanenti usati in pale eoliche;
    \item elementi utilizzati come catalizzatori nell'industria chimica;
    \item fosfori per l'illuminazione nelle lampade fluorescenti, poiché in molte giurisdizioni politiche queste lampade devono essere raccolte per isolare il mercurio pericoloso al loro interno;
    \item magneti utilizzati nei dischi rigidi dei server.
\end{itemize}
In seguito all'improvviso aumento dei prezzi delle terre rare esportate dalla Cina tra il 2010 e il 2011, molti ricercatori hanno avviato studi per ricercare metodi più efficaci per il recupero e il riutilizzo degli REE.
Nel 2014 Allan Walton, presso l’Università di Birmingham, ha sviluppato un processo di riciclo dei magneti al neodimio che consente di recuperare 1 kg di materiale consumando solo il 12\% dell’energia necessaria per l’estrazione da materie prime, con una riduzione del 98\% della tossicità per l’uomo. La procedura prevede l'esposizione dei magneti all’idrogeno: il neodimio lo assorbe, consentendo successivamente di separare, tramite setacciatura, la polvere di neodimio-ferro-boro dal rivestimento in nichel.
Attualmente, sono in corso ricerche in Giappone, presso l'Università di Tohoku, nell'\textit{European Rare Earth (Magnet) Recycling Network} e nel \textit{Critical Materials Institute} negli Stati Uniti.

\chapter{L'ascesa delle GPU}
La produzione e diffusione di schede grafiche (GPU) iniziò alla fine degli anni '70 con la realizzazione di dispositivi hardware per visualizzare testo sullo schermo dei personal computer.
Tra gli anni '90 e il 2000 aziende quali NVIDIA e ATI (ora chiamata AMD) iniziarono a produrre GPU in grado di fare 3D rendering grazie ad un architettura basata su silicone e un numero maggiore di transistor. Fino al 2010 le schede grafiche erano utilizzate quasi esclusivamente per la loro capacità di rappresentare accuratamente oggetti 3D su uno schermo, in particolare nei videogiochi e applicazioni professionali. Il loro vantaggio rispetto alle CPU deriva dall'utilizzo di molte "core" per eseguire in parallelo piccoli compiti.

\section{Produzione e impatto ambientale}

I materiali grezzi utilizzati da NVIDIA, l'azienda leader nella produzione di GPU con una quota di mercato globale pari a circa l’80\%, includono alluminio, lamine di rame, fibre di vetro, gel di silicone, stagno, titanio e tungsteno. NVIDIA è consapevole e dichiara pubblicamente che «stagno, titanio e tungsteno sono "minerali da conflitto", legati a violenze e violazioni dei diritti umani»; molti di questi materiali sono importati da Cina e Indonesia.
Il consumo di elettricità da parte di NVIDIA è per la gestione delle proprie strutture è la fonte principale dei loro gas serra, con circa 68.000 kg di emissioni. L'energia totale richiesta è aumentata del 33\% dal 2019-2020, quasi metà di questa è impiegata nelle operazioni nei loro data center.

Oltre a utilizzare 1/3 di energia rinnovabile dal 2020, NVIDIA afferma di aver ridotto i loro flussi di rifiuti liquidi inviati a discariche del 78\%. Si stanno anche facendo sforzi per implementare misure di conservazione idriche. Tuttavia, il loro fornitore TSMC usa 157,000 tonnellate d'acqua minerale al giorno per produrre semiconduttori. Non utilizzano acqua riciclata  poiché rischierebbe di contaminare l'elettronica.

I sistemi basati su GPU hanno un elevato impatto ambientale: sebbene una singola unità consumi meno energia rispetto a una CPU, il numero elevato di GPU utilizzate annulla questo vantaggio. A ciò si aggiungono i costi energetici legati al raffreddamento e le emissioni di gas serra generate dal loro funzionamento

\section{Utilizzi}

\subsection{AI e machine learning}

La capacità delle GPU di eseguire più compiti in parallelo ha cambiato il modo in cui le applicazioni di analisi e processo dei dati sono strutturate. Fornisce ad algoritmi per AI e machine learning la velocità ed efficienza necessaria per eseguire calcoli complessi.
In alcune GPU moderne sono presenti core specializzati per l'elaborazione di operazioni sui tensori, fondamentali nel deep learning. Questi consentono di accelerare le moltiplicazioni tra matrici, migliorando sia l'efficienza di addestramento che l'inferenza delle reti neurali.
Lo sviluppo di GPU ottimizzate per l'AI è un trend destinato a crescere in futuro, visti i continui progressi che si stanno facendo su quel fronte.

\subsection{Mining di criptovalute}

Nell'ultimo decennio, una varietà di valute digitali chiamate criptovalute si sono diffuse nel mercato mondiale. L'estrazione, o "mining", di criptovalute consiste nel raccogliere un insieme di transazioni effettuate negli ultimi dieci minuti all'interno di un blocco e risolvere un complesso problema matematico associato ad esso. Il primo miner che riesce a risolvere il problema aggiunge il blocco alla blockchain (una specie di registro digitale che contiene informazioni in modo sicuro e verificabile), rendendo ufficiali le transazioni contenute, e riceve come ricompensa un premio in Bitcoin.

Per risolvere questi problemi inizialmente erano impiegati gruppi di CPU, per poi essere sostituite da GPU. La crescente popolarità delle criptovalute, e il conseguente consumo energetico richiesto per il mining produce enormi quantità di diossido di carbonio. Inoltre, il rapido sviluppo di GPU sempre più performanti per il mining rende i modelli precedenti rapidamente obsoleti, producendo rifiuti elettronici (E-Waste). Come già evidenziato, il riciclo dei materiali contenenti terre rare è complesso e poco efficiente; una gestione inadeguata di questi rifiuti elettronici comporta il rilascio di sostanze nocive nel suolo e contribuisce all'inquinamento atmosferico.

L'impronta idrica del mining di criptovalute è una questione ambientale sempre più rilevante. Per raffreddare i gruppi computer utilizzati nelle operazioni di mining, si impiegano grandi quantità di acqua. Solamente per l'estrazione di Bitcoin, la valuta più vecchia e popolare, si sono utilizzati 1,600 gigalitri d'acqua in un anno: una singola transazione consuma tanta acqua quanto quella contenuta in una piscina domestica. Il consumo elettrico annuale è paragonabile a quello del Portogallo, e viene prodotta una quantità di E-Waste simile a quella dei paesi bassi.

Visti gli alti costi idrici ed elettrici per alimentare questi computer, sono stati presi provvedimenti legali per limitare l'impatto ambientale: nel 2021, la Cina ha vietato completamente il mining di Bitcoin, nel 2022 lo stato di New York ha vietato centri di mining alimentati da energia fossile. In Canada, la ditta elettrica del Quebec ha aumentato i prezzi e limitato l'utilizzo massimo di energia destinata al mining. Nel 2022 la commissione europea ha invitato gli stati membri a ridurre il consumo di elettricità dei miner.

\subsection{Veicoli a guida autonoma}

La sicurezza delle auto a guida autonoma dipende dalla loro capacità di computazione e di fare "scelte" rapide. La capacità delle GPU di eseguire più compiti in parallelo ottimizza compiti fondamentali per un'automobile a guida autonoma quali identificazione di oggetti e analisi delle immagini. NVIDIA fornisce servizi per lo sviluppo di sistemi di guida autonoma basati su modelli AI addestrati su moltissimi scenari di guida, testati e validati in uno spazio virtuale, NVIDIA Omniverse.

Il problema del surriscaldamento delle GPU si presenta anche qua: una macchina con GPU a bordo dovrebbe disporre di un notevole sistema di raffreddamento. Queste due componenti, assieme a continui aggiornamenti software, contribuiscono all'elevato prezzo delle auto a guida autonoma.

\subsection{Data center}

Con la diffusione di applicazioni client-server e la richiesta di sistemi decentralizzati, i data center sono cresciuti da computer grandi come una stanza a strutture iper-scalari, strategicamente posizionate al servizio di utenti da tutto il mondo. 

All'interno dei data center, le GPU coprono un ruolo cruciale nella gestione di calcoli complessi, algoritmi di machine learning e supporto per simulazioni scientifiche altamente precise. L'importanza delle GPU nei data center è destinata a crescere assieme alla quantità e complessità di dati che le aziende raccolgono. Un report dell'Agenzia dell'\textit{International Energy Agency} afferma che i data center consumano circa l'1-1.5\% dell'energia mondiale. Questa cifra può solamente aumentare più noi continuiamo a fare affidamento su servizi digitali. Questo report era stato fatto prima del boom dell'intelligenza artificiale, e Goldman Sachs stima che una singola query a ChatGPT richieda quasi 10 volte l'elettricità di una ricerca su Google. Per questo prevede che la domanda di energia dei data center crescerà del 160\% entro il 2030.

Secondo un report di The Guardian\cite{guardian2024datacenters}, tra le aziende del gruppo GAFAM, i data center di Amazon risultano essere i maggiori consumatori di energia, richiedendo più del doppio dell'elettricità rispetto alla seconda classificata, Apple. Un singolo data center di Amazon attualmente consuma l'equivalente elettrico di 50.000 case. Per questo motivo le GAFAM in primo piano, ma anche altri data center sono sotto pressione dai governi locali per diventare carbon neutral e passare a energie rinnovabili.

\chapter{Consumo d'acqua nei data center}

Le GPU rappresentano solo una parte dei macchinari elettronici che devono rimanere costantemente attivi all'interno dei data center; a queste si aggiungono router, hard disk, server e altre apparecchiature di rete, tutti elementi che contribuiscono in modo significativo al consumo energetico complessivo.

Nonostante l'enorme consumo energetico di queste strutture, la quantità d'acqua utilizzata per il raffreddamento è in realtà inferiore se paragonata a quella richiesta per la produzione di elettricità. Il processo di produzione di elettricità si basa sul riscaldamento dell'acqua, che, trasformandosi in vapore, mette in movimento una turbina e genera energia elettrica.
Per questo, i data center più grandi, che spesso gestiscono autonomamente la propria produzione di energia, utilizzano l'acqua anche per il raffreddamento dei loro sistemi elettrici e informatici. Un sistema elettrico autonomo implica, infatti, una maggiore uptime, ad esempio nel caso in cui i sistemi di alimentazione pubblici vengano meno a causa di un blackout.

Per quanto riguarda il raffreddamento, i sistemi a ciclo diretto prelevano l'acqua da corpi idrici vicini; l'acqua, riscaldata dal vapore prodotto nel condensatore, viene poi restituita alla fonte originale, ma a una temperatura superiore. Questi sistemi sono particolarmente vulnerabili in periodi di siccità e caldo estremo, poiché l'acqua restituita può superare i limiti termici consentiti, con impatti negativi sugli ecosistemi locali.
Per ridurre tali impatti, alcuni data center hanno adottato sistemi di raffreddamento a ciclo chiuso, che riciclano le acque di scarico e piovane, riducendo così l'uso di acqua fresca del 50-70\%. In questo tipo di sistema, l'acqua di raffreddamento non viene restituita alla sorgente, ma convogliata in torri di raffreddamento, dove il calore viene disperso tramite evaporazione. Il resto dell'acqua viene quindi ricircolato attraverso i condensatori.

Alcuni data center vanno oltre e scelgono di localizzare le proprie strutture in aree strategiche per ottimizzare il raffreddamento: per esempio, zone montane dove l'aria è naturalmente più fresca, o addirittura sott'acqua, come nel caso di Microsoft in Scozia, dove l'acqua fredda dell'oceano viene utilizzata per raffreddare le apparecchiature.

Gli Stati Uniti sono il paese con il maggior numero di data center, con più di 5300 sempre in funzione. I quantitativi di acqua consumati dai singoli data center sono difficili da stimare, poiché non sono di dominio pubblico e spesso oggetto di discussione, dato l'impatto negativo che possono avere sulle comunità locali in cui vengono costruiti. Si ipotizza che un data center di piccole dimensioni consumi circa 26 milioni di litri d'acqua all'anno tra raffreddamento e produzione di elettricità, e che Google abbia utilizzato 5 miliardi di litri solamente per il raffreddamento delle sue strutture.

\section{Conseguenze dell'uso di acqua minerale e potabile}

Con l'espansione dei servizi online, accelerata dal lockdown dovuto alla pandemia di Covid-19, si è registrato un aumento significativo nella costruzione di data center, spesso senza considerare adeguatamente le conseguenze economiche e ambientali legate alla posizione geografica scelta.
Uno studio fatto da Landon Marston, professore della Virginia Tech, afferma che l'industria dei data center statunitense preleva il 90\% della sua acqua da bacini idrogeologici, e che il 20\% di questi siano sovrasfruttati. In certi casi si usa anche acqua non potabile: Google afferma che il 25\%  delle risorse idriche impiegate derivino da prodotti fognari trattati. In media queste alternative all'acqua potabile costituiscono solo il 5\% del consumo totale di un data center. 

Si consideri il caso della città di Mesa, in Arizona, dove nel maggio 2021 è stato approvato lo sviluppo di un data center di 800 milioni di dollari in una zona in alta allerta per la siccità, la più alta registrata negli ultimi 126 anni. Il vicesindaco Jenn Duff ha sostenuto che "[...] i data center sono un uso irresponsabile della nostra acqua". Duff fa parte del numero crescente di individui che esprimono preoccupazioni circa la diffusione dei data center sul territorio nazionale, soprattutto in aree già aride, dove l'energia solare ed eolica è ampiamente disponibile e a prezzi bassi. 
Simili preoccupazioni hanno portato ad un divieto temporaneo alla costruzione di nuovi data center nei Paesi bassi e alla promulgazione di leggi per maggiore trasparenza in Francia.

Le città dei paesi in via di sviluppo promuovono la costruzione di data center per accelerare la trasformazione digitale, offrendo incentivi alle grandi aziende tecnologiche in cambio di investimenti e accesso ai mercati emergenti. Questo avviene senza considerare le conseguenze per il territorio e per le comunità locali, che spesso subiscono impatti negativi.
In India, si stima che il data center di Bengaluru consumi 8 milioni di litri d'acqua al giorno, minacciando i fragili sistemi idrici di una città che ha già vissuto una grave siccità.
Inoltre, durante la peggiore siccità registrata negli ultimi 74 anni, il governo dell'Uruguay ha autorizzato la costruzione di data center da parte di Google, provocando proteste tra la popolazione locale.

\section{Impatto sugli ecosistemi}
Il fabbisogno idrico dei data center è destinato a crescere, soprattutto con l'espansione di siti dedicati all'addestramento di sistemi di intelligenza artificiale. Si prevede che, a breve, il consumo d'acqua legato all'industria dell'AI supererà quello di settori tradizionalmente ad alto consumo idrico, come l'allevamento di bestiame e l'industria tessile. Inoltre l'impronta di CO$_2$ è maggiore del settore aereo.

I già citati sistemi di raffreddamento a ciclo diretto, fortunatamente ormai in disuso, erano responsabili della morte di miliardi di pesci all'anno, intrappolando nei loro sistemi di assunzione anche uova e larve. L'inquinamento termico provoca un aumento del metabolismo della fauna locale, con conseguente incremento del loro fabbisogno alimentare. Questo, a sua volta, porta a una carenza di cibo, spingendo le specie a migrare verso habitat più sostenibili. Il riscaldamento dell'acqua provoca anche una riduzione dei livelli di ossigeno dissolto nell'acqua, il quale è fondamentale per la sopravvivenza di organismi in ambienti acquatici.

Questi problemi non sono limitati ai sistemi di raffreddamento dei data center, ma di tutti gli impianti geotermici.

\section{Alternative sostenibili}

I data center più grandi, in grado di controllare la propria produzione di elettricità, hanno già provato a mettere in atto soluzioni sostenibili. Oltre all'esempio già menzionato di Microsoft in Scozia, il data center di Google ad Hamina, in Finlandia, utilizza acqua marina per il raffreddamento sin dalla sua inaugurazione nel 2011. Anche Amazon, seppur con scarsi risultati, si sta impegnando per essere meno dipendente dall'acqua minerale.

\subsection{Reattori modulari e fusione}
Per ridurre il consumo idrico e garantire un uptime prossimo al 100\%, molte big tech stanno investendo in fonti energetiche alternative. Sebbene le energie rinnovabili come l'eolico e il solare offrano il vantaggio di zero emissioni, non sono in grado di assicurare una fornitura continua di elettricità 24 ore su 24, 7 giorni su 7. Per questo motivo, l'area che sta attirando il maggior interesse e investimento è quella della fusione nucleare.

Gli Stati Uniti hanno prodotto una minuscola quantità di energia tramite fissione nucleare nel 2022, al Lawrence Livermore National Ignition Facility. Da allora investimenti nella fissione sono saliti fino a 8 miliardi di dollari. Di particolare rilievo è l'investimento da 1 miliardo di dollari effettuato da Sam Altman, CEO di OpenAI, in Helion, azienda specializzata nella ricerca sulla fusione nucleare. Questa iniziativa è motivata dall'elevato fabbisogno energetico necessario per sostenere infrastrutture come quella di ChatGPT, il cui addestramento e funzionamento costante comporta costi energetici estremamente elevati.

Per decenni, a causa dei rischi alla sicurezza e regolamentazioni, il nucleare non è stato ben visto negli Stati Uniti. Tuttavia, Amazon ha firmato accordi con Energy Northwest per la costruzione di reattori modulari (SMR) in grado di produrre eventualmente abbastanza energia per sostenere 770.000 case. Anche Google ha annunciato un accordo con Kairos Power per produrre 500 megawatt di elettricità tramite SMR entro il 2035.
Ha suscitato grande clamore l'annuncio di Microsoft, che in collaborazione con Constellation Energy ha deciso di riattivare l'impianto nucleare di Three Mile Island, oggetto nel 1979 di un parziale meltdown che portò al rilascio di gas radioattivi nell'atmosfera.


\chapter*{Fonti}
\renewcommand{\bibsection}{}
\section*{Riferimenti bibliografici e citazioni}
\nocite{*}

\bibliography{refs}
\section*{Paper scientifici}
\subsection*{Capitolo 1}
\urlstyle{same}
\sloppy

\url{https://www.sciencedirect.com/science/article/abs/pii/S0959652612006932} \\
\url{https://www.sciencedirect.com/science/article/abs/pii/S0176268022001598} \\
\url{https://www.gssc.lt/wp-content/uploads/2025/02/v04_Boruta_Rare-earths_A4_EN.pdf} \\
\url{https://www.sciencedirect.com/science/article/abs/pii/S2452223617301256} \\
\url{https://www.sciencedirect.com/science/article/pii/S0301420725000613?via%3Dihub}

\subsection*{Capitolo 2}
\url{https://neptjournal.com/upload-images/(55)B-4152.pdf}

\subsection*{Capitolo 3}
\url{https://www.nrdc.org/sites/default/files/power-plant-cooling-IB.pdf}\\
\url{https://www.nature.com/articles/s41545-021-00101-w}

\section*{Siti web e articoli di giornale}
\subsection*{Capitolo 1}
\url{https://hir.harvard.edu/not-so-green-technology-the-complicated-legacy-of-rare-earth-mining/} \\
\url{https://spectrum.ieee.org/ukraine-rare-earth-minerals} \\
\url{https://www.bbc.com/news/articles/c9d5jwvw9nlo} \\
\url{https://www.foxnews.com/politics/why-trumps-push-frigid-greenland-about-icing-out-us-adversaries} \\
\url{https://www.theguardian.com/us-news/2025/feb/25/trump-zelenskyy-ukraine-minerals-deal} \\
\url{https://www.soci.org/chemistry-and-industry/cni-data/2016/2/recyling-rare-earths}

\subsection*{Capitolo 2}
\url{https://retropcparts.com/gpu-manufacturing-process/} \\
\url{https://www.designlife-cycle.com/nvidia-gpu} \\
\url{https://tecex.com/gpu-powerhouses-behind-ai/} \\
\url{https://www.perplexity.ai/page/the-gpu-shortage-explained-ori-7BswhHKvT_idmwUL0P845Q} \\
\url{https://telnyx.com/resources/gpu-architecture-ai} \\
\url{https://medium.com/@vikaskumar8449/understanding-the-resources-consumed-by-ai-models-and-their-environmental-impact-b4fd9b2168d7} \\
\url{https://www.soprasteria.com/insights/details/what-is-the-actual-environmental-cost-of-ai} \\
\url{https://semianalysis.com/2024/03/13/ai-datacenter-energy-dilemma-race/} \\
\url{https://www.techtarget.com/searchdatacenter/feature/How-the-rise-in-AI-impacts-data-centers-and-the-environment} \\
\url{https://en.wikipedia.org/wiki/Environmental_impact_of_bitcoin} \\
\url{https://www.investopedia.com/tech/gpu-cryptocurrency-mining/} \\
\url{https://www.nvidia.com/en-us/self-driving-cars/} \\
\url{https://www.icdrex.com/gpus-for-self-driving-cars-powering-the-automotive-industrys-potential/} \\
\url{https://www.intel.com/content/www/us/en/products/docs/discrete-gpus/data-center-gpu/what-is-data-center-gpu.html} \\
\url{https://www.theguardian.com/technology/2024/sep/15/data-center-gas-emissions-tech} \\
\url{https://www.trgdatacenters.com/resource/the-role-and-purpose-of-data-center-gpus/}

\subsection*{Capitolo 3}
\url{https://www.asce.org/publications-and-news/civil-engineering-source/civil-engineering-magazine/} \\
\url{https://dgtlinfra.com/data-center-cooling/} \\
\url{https://dgtlinfra.com/data-center-water-usage/} \\
\url{https://www.nbcnews.com/tech/internet/drought-stricken-communities-push-back-against-data-centers-n1271344} \\
\url{https://eng.ox.ac.uk/case-studies/the-true-cost-of-water-guzzling-data-centres/} \\
\url{https://www.lawfaremedia.org/article/ai-data-centers-threaten-global-water-security} \\
\url{https://www.weforum.org/stories/2024/11/circular-water-solutions-sustainable-data-centres/}

\newpage
\end{document}
